% AY2021 制御工学 第2問〔II〕（転記）
% 出典: 03_制御工学/original/03_制御工学/2021/2021se.pdf
% 要約: G1,G2,G3 の直列に2つの負帰還（G2 まわり、および G1 後段ノードから入力側）。
%       伝達関数・ゲイン線図・ステップ応答・G3 変更時の過渡特性比較。構造は AY2016 第2問と同型。

\section*{第2問〔II〕}

制御システムに関する次の問い (1)〜(4) に答えよ.

\begin{center}
\begin{tikzpicture}[>=Latex]
  \node (in)   at (0,0)   {$U(s)$};
  \node (sum1) [sum] at (1.4,0) {};
  \node (G1)   [block] at (2.9,0) {$G_1(s)$};
  \node (sum2) [sum] at (4.6,0) {};
  \coordinate (nA) at (5.4,0);
  \node (G2)   [block] at (6.7,0) {$G_2(s)$};
  \coordinate (nB) at (8.0,0);
  \node (G3)   [block] at (9.2,0) {$G_3(s)$};
  \node (out)  at (10.8,0) {$Y(s)$};

  \draw[->] (in)   -- (sum1);
  \draw[->] (sum1) -- (G1);
  \draw[->] (G1)   -- (sum2);
  \draw[->] (sum2) -- (G2);
  \draw[->] (G2)   -- (G3);
  \draw[->] (G3)   -- (out);
  \fill (nA) circle (1.4pt);
  \fill (nB) circle (1.4pt);

  \draw[->] (nB) -- (8.0,1.3) -- (4.6,1.3) -- (sum2.north);
  \draw[->] (nA) -- (5.4,-1.3) -- (1.4,-1.3) -- (sum1.south);

  \node at ($(sum1.north west)+(0.02,0.10)$) {\scriptsize $+$};
  \node at ($(sum1.south east)+(-0.02,-0.10)$) {\scriptsize $-$};
  \node at ($(sum2.north east)+(-0.02,0.10)$) {\scriptsize $-$};
  \node at ($(sum2.south west)+(0.02,-0.10)$) {\scriptsize $+$};
\end{tikzpicture}
\end{center}

\begin{enumerate}[label=(\arabic*)]
  \item 下図で表されるシステムの伝達関数を求めよ.

  \item 問い (1) で求めた伝達関数に, $G_1(s)=\dfrac{3}{s}$, $G_2(s)=\dfrac{1}{s+4}$,
        $G_3(s)=4$ を代入してゲイン (dB) を求め, ゲイン線図を漸近線によって描け.
        折れ点等の特徴点を明記し, 解答用紙内の縦軸・横軸の目盛には適宜数値を記入すること.

  \item 問い (2) で求めたシステムのステップ応答を求め, 時間応答の概形を描け.

  \item $G_3(s)$ を $2s+4$ に変更した場合のステップ応答を求めよ. また, このシステムと
        問い (2) で求めたシステムに関して, ゲイン線図ならびに時間応答の両観点から比較し,
        過渡特性の違いについて考察せよ.
\end{enumerate}
